\subsubsection{TempoPFN synthetic prior used by Toto 2.0}

\paragraph{Terminological distinction.}
Two similarly named works must be separated. \emph{TempoPFN}, by
\citet{moroshan2025tempopfn}, is the synthetic-data method cited and used by
Toto 2.0 \citep{khwaja2026toto2}. \emph{TimePFN}, by
\citet{taga2025timepfn}, is a different multivariate forecasting model based on
LMC-Synth; it is analyzed in the next subsection. Toto 2.0 consumes data from
the TempoPFN prior, rather than the TimePFN/LMC-Synth corpus.

\paragraph{How the TempoPFN prior generates series.}
TempoPFN defines a portfolio of complementary generator families rather than
one equation family \citep{moroshan2025tempopfn}. Adapted generators include
ForecastPFN component compositions, KernelSynth, a broader GP prior, and
CauKer structural causal models. New procedural families add sawtooth waves,
piecewise-constant steps, anomaly and spike processes, configurable sine
mixtures, and audio-inspired stochastic rhythms, financial volatility,
network traffic, and multiscale fractal signals. Its most structured stochastic
generator is a regime-switching, time-inhomogeneous Ornstein--Uhlenbeck process,

\begin{equation}
 d y_t
 =\theta(t,r_t)\bigl(\mu(t,r_t)-y_t\bigr)\,dt
  +\sigma(t,r_t)\,dW_t,
\label{eq:tempopfn-ou}
\end{equation}

where the latent regime $r_t\in\{0,1\}$ follows a Markov chain. The mean-
reversion speed, time-varying mean, and volatility can change between regimes
and evolve through polynomial, sinusoidal, logistic, or piecewise-linear
trends. Seasonal effects can modify the mean and volatility, and fractional
Brownian motion can replace $W_t$ to introduce long memory.

After a base trajectory is generated, an offline augmentation cascade expands
the prior. It optionally normalizes the source, applies an early convex
TS-Mixup, samples transformations from invariance, structure, seasonality,
signal-processing, discrete-effect, and measurement-artifact categories,
optionally applies random 1D convolution filters, performs a late TS-Mixup,
and finishes with scaling and low-amplitude noise. Examples include time
reversal, sign inversion, change points, shocks with exponential recovery,
calendar effects, amplitude modulation, differentiation or integration,
censoring, quantization, and resampling artifacts. An online stage injects
missing-value patterns. Consequently, one stored trajectory may combine a
sampled base mechanism with several independently recorded transformations.

\paragraph{High-throughput data construction.}
TempoPFN reports a corpus of approximately 10 million series with lengths up to
2048. Most procedural generators run on CPU, while neural generators such as
CauKer use a GPU. On the authors' dual 32-core AMD EPYC system, throughput for
length-2048 series ranges from $0.32$ series/s for exact KernelSynth and $0.66$
for CauKer to more than $100$ series/s for step, sine, anomaly, spike, and
sawtooth generators. This heterogeneous architecture is deliberate: expensive
GP/causal priors contribute structural richness, while inexpensive procedural
priors supply volume. During training, the loader samples a total length from
$\{128,256,512,1024,1536,2048\}$, favors length 2048, shortens by cutting or
subsampling, and then samples a history--future split. Mixed batches therefore
cover several mechanisms, context lengths, and forecast horizons.

\paragraph{How Toto 2.0 uses this prior.}
The Toto 2.0 report states that its base models are pretrained on 42.5\%
Datadog observability data and 57.5\% TempoPFN synthetic data
\citep{khwaja2026toto2}. The three larger models see 5.04 trillion points, of
which 2.90 trillion are synthetic; the 4M- and 22M-parameter models see 3.40
trillion points with the same proportions. The report identifies the
TempoPFN method and its coverage of nonstationary trends, abrupt change points,
and long-range dependencies. It reports the final mixture rather than an
itemized count for each TempoPFN generator inside the Toto corpus.

Toto 2.0 uses a decoder-only patched Transformer with patch size
$P_{\mathrm{patch}}=32$. Transformer blocks alternate causal attention along
time with full attention across variates. Contiguous patch masking replaces
autoregressive patch generation: during training, the model reconstructs
random contiguous masked spans; during inference, mask tokens occupy the
complete future horizon and the Transformer predicts that horizon in one
forward pass. A block-decoding mode remains available for very long horizons.
The output head predicts nine quantiles and is trained with pinball loss.

Three engineering choices support large-scale training. First, robust causal
$\operatorname{asinh}$ normalization controls the extreme scales and spikes
found in observability metrics. Second, NorMuon optimizes internal matrix
parameters while AdamW handles projections, biases, and normalization
parameters. Third, unit-scaled maximal-update parametrization (u-$\mu$P) lets
the authors tune a 10M-parameter proxy and transfer its learning configuration
to 4M, 22M, 313M, 1B, and 2.5B variants. Their released
\texttt{dd\_unit\_scaling} layer preserves this parametrization through
\texttt{torch.compile}, FSDP2 sharding, and data/tensor parallelism. All model
sizes train with 4096-step contexts, 32 variates per sample, and global batch
size 64. The generator portfolio provides data diversity; patch masking,
alternating attention, u-$\mu$P, and distributed sharding provide model-
training scale.

\paragraph{Explanation information available from the TempoPFN pipeline.}
The generator identity, sampled parameters, regime path, and augmentation
sequence can form a hierarchical mechanism record. For example,
Eq.~\eqref{eq:tempopfn-ou} supplies regime- and time-dependent drift and
diffusion parameters, while a step generator supplies exact change-point
locations. Their current use in Toto 2.0 is forecasting pretraining. For the
present thesis, retaining this metadata would permit evaluation by mechanism
family and controlled difficulty, whereas local gradient or Shapley ground
truth would be calculated from an explicitly selected executable forecasting
functional, such as the conditional mean of the sampled process.

